\section{Ejercicios}
\subsection{Hoja 1}
\begin{problem}[5]
Dada una sucesión $\lbrace f_n \rbrace \in R([a,b])$ que converge uniformemente a $f$, se pide demostrar que $f$ es ingrable Riemann y que:
\[ lim \int_{a}^{b} f_n = \int_{a}^{b} f \]

\solution
Supongamos que f es integrable Riemann, entonces tenemos que ver que:
\[ \forall \epsilon > 0 \ , \exists N \tq \forall n> N \ \abs{\int_a^b f_n - \int_a^b f} < \epsilon\]

Sabemos que:
\[\abs{\int_a^b f_n - \int_a^b f} \leq \abs{\int_a^b \abs{f_n -f}} dx\]

Recordemos la definición de convergencia uniforme

\begin{defn}[Convergencia\IS uniforme]
\[f_n \xrightarrow{uniforme} f \Leftrightarrow \forall \epsilon < 0, \exists N_{\epsilon} \tq \forall x \in [a,b]  \forall n \geq N_{\epsilon}, \abs{f_n (x) - f(x)} < \epsilon\]
\end{defn}

Si $n \geq N_{\frac{\epsilon}{b-a}}$ entonces, usando la definición de convergencia uniforme:

\[\int_a^b \abs{f_n(x) - f(x) dx} \leq \int_a^b \frac{\epsilon}{b-a}dx = \epsilon\]

Por tanto queda claro que si $f$ es integrable Riemann podemos conmutar el límite con la integral. Ahora queda ver por qué $f$ es integrable Riemann.

$f$ será integrable Riemann sii:
\[\forall \epsilon > 0 \ \exists P \tq \forall P' \prec P \]
\[\overline{J}_{P'}(f) - \underline{J}_{P'}(f) < \epsilon\]

Vamos a probarlo:
\[\overline{J}_P(f) - \underline{J}_P(f) = \sum(\sup_k f - \underset{k}{inf} f)\abs{I_k} \leq\]
\[\leq \sum\left( \abs{\sup_{k}f_n(x) - \sup_{k}f(x)} + \abs{\sup_{k}f_n(x) - \underset{k}{inf}f_n(x)} +  \abs{\underset{k}{inf}f_n(x) - \underset{k}{inf}f_n(x)} \right)\abs{I_k} \leq \]
Puesto que $f_n$ converge uniformemente a $f$ habrá un $n$ a partir del cual la distancia máxima entre $f_n$ y $f$ sea $\frac{\epsilon}{6}$ y por tanto la distancia máxima entre el supremo y el ínfimo de $f_n$ será menor que $\frac{\epsilon}{3}$
\[\leq \sum\left( \sup_{k} \abs{f_n(x) - f(x)} + \frac{\epsilon}{3} + \sup_k \abs{f_n(x) - f(x)} \right)\abs{I_k} \leq\]
Aplicando de nuevo la convergencia uniforme y tomando el máximo entre este $n$ y el calculado en el paso interior nos queda:
\[\leq \sum \left(\frac{\epsilon}{3} + \frac{\epsilon}{3} + \frac{\epsilon}{3} \right)\abs{I_k} = \sum\epsilon\abs{I_k}\]

Como hay convergencia uniforme entre $f_n$ y $f$ podemos hacer que los supremos sean tan pequeños como queramos y hacer así que el interior sumatorio quede menor que $\epsilon \abs{I_k}$

Puesto que $\epsilon$ es un número cualquiera podemos hacerlo tan pequeño como queramos haciendo que el último sumatorio escrito tienda a 0.

\end{problem}

\begin{problem}[6]

Sea $\lbrace f_n \rbrace$ una sucesión monótona creciente de funciones continuas en un intervalo $I=[a,b]$ que convergen en dicho intervalo a otra función continua $f$. Demuestra que entonces:
\[ lim \int_{a}^{b} f_n(x) = \int_{a}^{b} f(x) \]
\solution
Por ser funciones continuas son integrables Riemann. Si conseguimos demostrar que convergen uniformemente podemos emplear el ejercicio anterior y lo tendríamos hecho.
\end{problem}

\begin{problem}[7]
Dada la sucesión $I_k = (a_k, b_k)$ tales que $\bigcup_{k=1}^{N}~I_k~=~[a,b]$

Demostrar que:
\[b-a \leq \sum_{k=1}^N (b_k - a_k)\]

\solution
Vamos a utilizar la integral de Riemann como recomienda el ejercicio, utilizando la función indicatriz de cada intervalo $\ind_{I_k}$

Está claro que la función indicatriz del intervalo I es menor o igual que la suma de las funciones indicatrices de los intervalos. Lo cual es obvio, ya que si $x$ está en el intervalo, $x$ estará también en al menos uno de los intervalos $I_k$. Es decir:
\[\ind_{[a,b]} \leq \sum_{k=1}^{N} \ind_{I_k}\]

Utilizando la monotoreidad de la integral de Riemann podemos ``integrar a ambos lados`` obteniendo:

\[\int_{a}^{b} \ind_{[a,b]} \leq \int_{a_k}^{b_k} \sum_{k=1}^{N} \ind_{I_k}\]

Por la linealidad de la integral, podemos incluso meter la integral denro del sumatorio
\[\int_{a}^{b} \ind_{[a,b]} \leq \sum_{k=1}^{N} \int_{a_k}^{b_k} \ind_{I_k}\]


Conociendo la integral de la función indicatriz tenemos el resulado de forma inmediata.
\[b-a \leq \sum_{k=1}^{N} ( b_k - a_k )\]

\end{problem}

\begin{problem}[9]
Dado $O\subset (a,b)$ unión numerable de intervalos disjuntos, se pide demostrar:
\[m^*(O)=\sum_{n=1}^{\infty} \abs{I_k}\]

\solution
Está claro que:
\[m^*(O) \leq \sum_{n=1}^{\infty} \abs{I_k}\]

Tenemos que demostrar la desigualdad contraria para poder concluir la igualdad.
Vamos a probar que:
\[m^*(O) \geq \sum_{n=1}^{\infty} \abs{I_k} - \epsilon\]

Puesto que sabemos que la suma infinita tiene un resultad finito (ya que $O$ es finito)
\[\forall \epsilon > 0 \exists N \tq \sum_{n=1}^N\abs{I_n} \geq \sum_{n=1}^{\infty} \abs{I_k} - \frac{\epsilon}{2}\]

Vamos a definir los intervalos $I'_n$ que serán un poco más pequeños que los $I_n$
\[\forall n=1,2,...,N \text{ definimos } I'_n=[a_n+\frac{\frac{\epsilon}{2}}{2^{n+1}}, b-\frac{\frac{\epsilon}{2}}{2^{n+1}}]\]

Ahora tenemos que:
\[\sum_{n=1}^N \abs{I'_n}=\sum_{n=1}^N\abs{I_n} - \frac{\epsilon}{2}\sum_{n=1}^{N}\frac{1}{2^n} \geq \sum_{n=1}^N\abs{I_n} - \frac{\epsilon}{2}\]

Definimos ahora el compacto:
\[K_n = \bigcup_{n=1}^N I'_n\]

Sean $J_n$ intervalos abiertos tales que $O \subset \bigcup_{n=1}^{\infty}J_n \Rightarrow K_n \subset \bigcup_{n=1}^{\infty}J_n$.

Por tanto
\[\exists N' \tq K_n \subset \bigcup_{n=1}^{N'}J_n=A_N\]

Vamos a fijarnos ahora en la medida exterior de $O$:

\[ m^*(O) = \inf\lbrace \sum_{n=1}^{\infty}\abs{J_n} \rbrace \geq \inf\lbrace \sum_{n=1}^{N'}\abs{J_n} \rbrace \geq \]
\[ \geq \inf\lbrace \int \ind_{A_N} \rbrace \geq \inf\lbrace \int \ind_{K_N} \rbrace \geq \inf\lbrace \sum_{n=1}^N\abs{I'_n} \rbrace \geq \]
\[ \geq \inf\lbrace \sum_{n=1}^N\abs{I_n} -\frac{\epsilon}{2}\rbrace \geq \inf\lbrace \sum_{n=1}^N\abs{I_n} -\epsilon \rbrace \]

Y obtenemos así la desigualdad buscada.
\end{problem}

\begin{problem}[10]
Dado un compacto K contenido en el intervalo (a,b), nos piden demostrar que:
\[m(K)<b-a\]

\solution
Consideremos la sucesión: $(a+\frac{1}{n}, b - \frac{1}{n})$, con $n >\frac{2}{b-a}$. Obviamente:
\[K \subset \bigcup_{n=n_0}^{\infty}(a+\frac{1}{n}, b - \frac{1}{n})\]

Como K es compacto:
\[\exists N \tq K \subset \bigcup_{n=n_0}^{N}(a+\frac{1}{n}, b - \frac{1}{n}) = (a+\frac{1}{N}, b - \frac{1}{N}) \]

Y llegamos a:
\[m(K) \leq b-a-\frac{2}{N} < b-a\]

\end{problem}

\begin{problem}[11]
Sea $C_n$ una sucesión creciente de subconjuntos medibles contenidos en (a,b) y sea $C$ la unión de estos subconjuntos. Se pide demostrar que:
\[m(C_n) \rightarrow m(C)\]

\solution
Vamos a definir la sucesión $D_n$ como:
\[D_1=C_1 \ D_2 = C_2 \setminus D_1 \ D_3 = C_3 \setminus (C_1 \bigcup C_2) \ ...\]

Teniendo así C expresado como unión de los $D_n$, que son disjuntos. Así, la medida de $C$ queda expresada como:
\[m(C)=\sum_{n=1}^{\infty}m(D_n)\]
sabiendo que:
\[m(C_N)=\sum_{n=1}^{N}m(D_n)\]
\end{problem}

\begin{problem}[12]
Sea $C_n$ una sucesión decreciente de subconjuntos medibles contenidos en (a,b) y sea $C$ la intersección de estos subconjuntos. Se pide demostrar que:
\[m(C) \rightarrow m(C)\]
\solution

Tomando los complementarios, que también son medibles, tenemos una sucesión creciente de subconjuntos medibles contenidos en (a,b). Aplicando el ejercicio anterior llegamos a que:
\[(b-a)-m(C_n) \rightarrow (b-a)-m(C)\]
De donde puede deducirse que $m(C_n)$ decrece hacia $m(c)$.
\end{problem}

\begin{problem}[14]
Dada una sucesión $A_n$ contenia en (a,b), se pide demostrar que:
\[\lim_n m(A_0\Delta A_n)=0 \Rightarrow \lim_n m(A_n)=m(A_0) \]

Recordemos que:
\[ A_n \Delta A_0 = (A_n \setminus A_0) \cup (A_0 \setminus A_n) =
(A_n \cup A_0)\cap(A_n^c \cup A_0^c) \]
\solution
\[\lim_n m(A_0\Delta A_n)=0 \Rightarrow \lim_n m(A_0^c \cap A_n)=0 \wedge \lim_n m(A_0\cap A_n^c)=0\]

Escribimos ahora $A_n$ y $A_0$ como:
\[A_n = (A_n \cap A_0) \bigcup (A_n \cap A_0^c)\]
\[A_0 = (A_0 \cap A_n) \bigcup (A_0 \cap A_n^c)\]

De aquí podemos ver que:
\[m(A_n) = m(A_n \cap A_0) + m(A_n \cap A_n^c)\]
\[m(A_0) = m(A_0 \cap A_n) + m(A_0 \cap A_n^c)\]

Y restando llegamos a:
\[m(An) = m(A_0) -m(A_0 \cap A_n^c) + m(A_n \cap A_0^c)\]

Y sabemos que las medidas de estas intersecciones tienden a 0.
\end{problem}

\subsection{Hoja 2}

\begin{problem}[2]
Sea $X=\{a,b,c,d\}$ y sea $\algb{E}=\{\{a\},\{b\}\}$. Se pide construir la $\salgb$ generada por $\algb{E}$

\solution
Sabemos que en un conjunto finito toda álgebra es una $\salgb$, ya que la diferencia entre álgebra y $\salgb$ era el cierre por uniones infinitas numerables. Si un conjunto es finito, $\algb{P}(X)$ será finito y no habrá posibilidad de hacer uniones infinitas.

Vamos a construir la mínima álgebra que contiene a $\algb{E}$ a pelo, forzando que se cumplan las propiedades de un álgebra de conjuntos:
\[\algb{M}(\algb{E})=\{\emptyset, X, \{a\}, \{b\}, \{a,b\}, \{b,c,d\}, \{a,c,d\}, \{c,d\}\}\]

Con este ejemplo podemos comprobar sencillamente que:
\[\algb{M}(A_1 \cup A_2) \neq \algb{M}(A_1) \cup \algb{M}(A_2)\]
Tomando $A_1=\{a\} \ y \ A_2 = \{b\}$

\end{problem}

\begin{problem}[4]
Dada una función $\appl{g}{X}{Y}$ y sea $\algb{A}$ una $\salgb$ de X, demostrar que:
\[\algb{B}=\{E\subset Y \tq g^{-1}(E)\in \algb{A}\}\]
es una $\salgb$

\solution
Basta con ver que se cumplen las siguientes dos propiedades
\begin{enumerate}
\item 
\[g^{-1}(E_1 \cup E_2)=g^{-1}(E_1) \cup g^{-1}(E_2)\]
\item
\[g^{-1}(Y \setminus E) = X - g^{-1}(E)\]
\end{enumerate}


\begin{proof}
\begin{enumerate}
\item
\[x\in g^{-1}(E_1 \cup E_2) \iff g(x) \in E_1 \cup E_2 \iff  \exists i=1,2 \ g(x) \in E_i \iff \]
\[\iff \exists i=1,2 \ x\in g^{-1}(E_i) \iff x\in g^{-1}(E_1) \cup g^{-1}(E_2) \]
\item
\[x\in g^{-1}(Y \setminus E) \iff g(x) \in Y \setminus E \iff g(x) \notin E \iff x \notin g^{-1}(E) \iff x \in X\setminus g^{-1}(E)\]
\end{enumerate}

\end{proof}
Así, por la primera propiedad queda claro que $\algb{B}$ es cerrado por uniones y con la segunda propiedad vemos que es cerrado por complementación. Es decir:
\[\{E_i\}_{i_1}^{\infty} \in \algb{B} \implies \bigcup_{i=1}^{\infty} E_i \in \algb{B}\]
ya que $g^{-1}(\bigcup_{i=1}^{\infty} E_i) = \bigcup_{i=1}^{\infty} g^{-1}(E_i) \in A$

Podemos hacer lo mismo para ver el cierre por complementación

Además, queda claro que tanto el vacío como el total se contienen en $\algb{B}$, ya que $X$ y $\emptyset \in A \implies g^{-1}(Y)=X \in A, g^{-1}(\emptyset)=\emptyset \in A \implies Y, \emptyset \in \algb{B}$ por lo que se trata de una $\salgb$
\end{problem}

\begin{problem}[5]
Dada una función $\appl{g}{X}{Y}$ y sea $\algb{B}$ una $\salgb$ de Y, demostrar que:
\[\algb{A}=\{g^{-1}(E) \tq E \in \algb{B}\}\]
es una $\salgb$

\solution
Vamos a comprobar las propiedades de cierre por unión y por complementación:
\begin{enumerate}
\item. Debemos ver que:
\[A_1, A_2\in \algb{A} \Rightarrow A_1 \cup A_2 \in A\]
Lo cual es cierto ya que, como demostramos en el ejercicio anterior:
\[g^{-1}(E_1 \cup E_2)=g^{-1}(E_1) \cup g^{-1}(E_2)\]
Si tomamos $A_i=g^{-1}(E_i)$ queda claro que:
\[A_1 \cup A_2 = g^{-1}(E_1) \cup g^{-1}(E_2) = g^{-1}(E_1\cup E_2) \Rightarrow g^{-1}(E)\]
para algún $E\in \algb{B}$ por ser una $\salgb$.

\item Ahora tenemos que ver que:
\[A \in \algb{A} \Rightarrow X\setminus A \in \algb{A}\]
Con la segunda propiedad del ejercicio anterior queda claro que es cierto
\end{enumerate}

Por tanto, $\algb{A}$ cumple todas las propiedades de $\salgb$.
\end{problem}

\begin{problem}[0-0]
\textbf{Inventado por el profesor.}

Vamos a ver que la imagen directa de una $\salgb$ no tiene por qué ser una $\salgb$, es decir:
dada una función $\appl{f}{X}{Y}$ y sea $\algb{A}$ una $\salgb$ de X, demostrar que:
\[\algb{B}=\{g(E) \tq E \in \algb{A}\}\]
no es necesariamente una $\salgb$

\solution
Esto es sencillo puesto que si la función $f$ no es suprayectiva, entonces $Y \neq g(E)$ para cualquier $E$, por lo que $\algb{B}$ no contiene al total y, por tanto no sería $\salgb$.

Supongamos ahora que es suprayectiva la función $f$. En este caso seguríamos teniendo problemas si $f$ no es inyectiva. 

Tomemos por ejemplo $g(x)=x^2$. En este caso, $g((-\infty, 0] \cup [0, \infty))=g(0)=0 \neq g((-\infty, 0]) \cup g([0, \infty))$.
\end{problem}

\begin{problem}[6]
Demuestra que una álgebra $\algb{A}$ es una $\salgb$ si y sólo si es cerrada para uniones numerables crecientes.

\solution
Vamos a demostrar las dos direcciones de la implicación:
\begin{itemize}
\item $\Rightarrow$
Es obvio ya que una $\salgb$ es un álgebra cerrada por uniones numerables y por tanto, es cerrada para el caso concreto de uniones numerables crecientes.
\item $\Leftarrow$
\textcolor{red}{Hecho por mi. No es seguro al 100\%}

Tenemos que demostrar que dada una unión numerable cualquiera $\{A_i\}_{i=1}^{\infty} \in \algb{A}$, se cumple que $\bigcup_{n=1}^{\infty} A_i \in \algb{A}$. Vamos a ello.

Dada mi unión numerable, construyo otra de la siguiente forma:
\[B_n = \bigcup_{i=1}^{n} A_i\]
de tal forma que $B_i \subset B_j \forall i<j$.

Ahora, sabiendo que la unión de los $B_i$ se contiene en la $\salgb$, tenemos que:
\[\bigcup_{n=1}^{\infty} A_i=\bigcup_{n=1}^{\infty} B_i \in \algb{A}\]

\end{itemize}
\end{problem}

\begin{problem}[7]
Determina el álgebra $\algb{A}$ generada por la colección de los subconjuntos finitos de un conjunto X no-numerable. Determina la $\salgb$ generada por $\algb{A}$. Estudiar el mismo problema en caso de que el conjunto X sea infinito numerable

\solution
Dado el conjunto de todos los subconjuntos finitos, para convertirlo en una $\salgb$ debemos asegurarnos de que sea cerrado por uniones numerables y por complementación.

Para hacerlo cerrado por uniones numerables, debemos incluir todos los conjuntos numerables, ya que cualquiera de estos puede obtenerse como unión de conjuntos finitos.

Además, si queremos que sea cerrado por complementación, tendremos que incluir los complementarios de todos los conjuntos mencionados anteriormente, es decir, debemos incluir todos aquellos conjuntos cuyo complementario sea numerable.

Así, vamos a construir de forma directa la $\salgb$ pedida:
\[\algb{M}(\algb{E})=\{ E\subset X \tq E \text{ numerable o finito, } E^c \text{ numerable o finito}\}\]

Podemos comprobar fácilmente que es una $\salgb$ observando que cumple las propiedades necesarias siguiendo el mismo razonamiento que el realizado para construirla. (está hecho en teroía). Es la mínima por construcción.

Si X es infinito numerable entonces 
\[\algb{M}(\algb{E})=\algb{P}(X)\]
aplicando la misma regla de construcción, ya que todos los subconjuntos de $\algb{P}(X)$ son numerables.
\end{problem}

\begin{problem}[8]
La $\salgb$ de (0,1] engendrada por:
\[\algb{E}= \{(0, \frac{1}{n}]: n=1,2,...\}\]
está formada por uniones finitas o numerables de intervalos (a,b]. Estudia cómo son estos intervalos.

\solution
Queremos ver como son los intervalos (a,b] tales que:
\[\algb{M}(\algb{E})=\{\bigcup_{i=1}^{\infty}(a_i, b_i]\}\]

Vamos a ver cuáles son los elementos que tenemos en $\algb{M}(\algb{E})$. Aquí, además de los propios elementos de $\algb{E}$ tenemos:
\begin{itemize}
\item \textbf{Complementarios} Son de la forma $(\frac{1}{n}, 1]$

\item \textbf{Uniones} Son de la forma $(\frac{1}{m}, \frac{1}{n}]$ con m>n

\item \textbf{Intersecciones} Intersecando dos elementos seguimos estando en $\algb{E}$, así que no ganamos nada nuevo.
\end{itemize}

Puesto que las uniones contienen a los complementarios, tenemos que los intervalos que forman la $\salgb$ son de la forma:
\[\algb{M}(\algb{E})=\{\bigcup_{i=1}^{\infty}(a_i, b_i]: a_i =\frac{1}{m}, \ b_i = \frac{1}{n} \ m>n\}\]
\end{problem}

\begin{problem}[9]
Describe la $\salgb$ generada por:
\[\algb{E}= \{N \subset \nat \tq \forall n \in \nat, \ 2n \in N\}\]

\solution
Vemos que el conjunto $\algb{E}$ está formado por todos los conjuntos que contienen a todos los pares.

Esta claro que la unión y la intersección de conjuntos de $\algb{E}$ pertenecen a $\algb{E}$ ya que contendrán a todos los pares.

Los complementarios serán los subconjuntos de $\nat$ que sólo contengan números impares.

Si combinamos todos estos conjuntos formaremos la mínima $\salgb$ que lo contenga todo. Es decir:
\[\algb{M}(\algb{E})=\{N \subset \nat \tq \text{Todos los n son impares, o N contiene a todos los pares}\}\]
\end{problem}

\begin{problem}[10]
Sea $\algb{M}$ una $\salgb$ de cardinal infinito. Demuestra que tiene cardinal no numerable.
\solution
Si la $\salgb$ es infinita podemos encontrar una colección numerable y disjunta de elementos de la misma. (Lo podemos hacer tomando una colección numerable y haciéndola disjunta mediante la eliminación en cada elemento de la unión de los anteriores).

Denotamos a esta colección como: $\{A_n:n\in \nat\}$. 

Ahora vamos a construir una colección infinita de $A_x$ como sigue:
\[\forall x \in (0,1) \text{ obtenemos el desarrollo en base 2 de } x=0,x_1,x_2,x_3,...\]
\[A_x=\bigcup_{n=1}^{\infty}A'_n \text{ donde } A'_n=\left\{ \begin{array}{lcc}
             A_n &   si  & x_n = 1 \\
             \\ \emptyset &  si  & x_n \neq 1
             \end{array}
   \right.\]

Como no puede darse el caso de dos $A_x$ iguales, es necesario tomarlos todos para construir $\algb{M}$ y puesto que el número de $A_x$ es no numerable (uno por cada real en el intervalo (0,1)), tenemos que el número de elementos de $\algb{M}$ es no numerable.

\end{problem}

\begin{problem}[11]
Hallar una cota superior al número de elementos que puede tener una $\salgb$ $\algb{M}$ generada a partir de un conjunto con n elementos:
\solution
El enunciado nos dice que tenemos un conjunto $ε\subset \algb{P}(X) \tq ε=\{A_1,...A_n\}$

Si el conjunto ε es una partición del conjutno $X$ (división de $X$ en subconjuntos disjuntos dos a dos), entonces la mínima $\salgb$ generada tendrá $2^{\#ε}$ elementos. 
\begin{proof}
Vamos a considerar todas las posibles uniones de los elementos de ε, que deberán pertenecer a $\algb{M}(ε)$.
Para poder contar estas uniones, vamos a representar cada unión mediante un número binario de longitud \# ε, donde los 1s nos indican qué elementos de ε estamos considerando.

Queda claro así que el conjunto de todas las uniones posibles tiene $2^{\# ε}$ elementos. Estas uniones serán todas disjuntas puesto que así lo son los elementos que las generan. Además, esto implica que las intersecciones son vacías y el complementario de cada elemento está formado por la unión del resto y, por lo tanto, también se contiene en el conjunto de uniones. 

Obviamente esta construcción genera la mínima $\salgb$ generada por el conjunto ε y que tiene el cardinal buscado.
\end{proof}

Pero el caso que nos concierne es ligeramente distinto, pues no son disjuntos.

Si tomamos el primer elemento, podemos construir una partición de $X$ a partir de él, considerando los conjuntos: $A_1 \ A_1^c$.

Cogiendo ahora un segundo conjunto $A_2 \in ε$, tenemos entonces que los cuatro conjuntos: $A_1 \cap A_2, \ A_1^c \cap A_2 \ A_1^c \cap A_2^c \ A_1 \cap A_2^c$, se contendrían en la $\salgb$. Estos cuatro conjuntos podrían ser distintos todos ellos, o podrían coincidir. Como mucho darán cuatro conjuntos distintos y como poco 2.

Para dejarlo claro, los hemos conseguido considerando las posibles intersecciones del segundo conjunto que hemos cogido y su complementario con los conjuntos que teníamos en el paso anterior.

Tomo ahora un tercer elemento de ε distinto de los anteriores, y construyo todas las combinaciones posibles con los 4 elementos anteriores, de forma similar a como construimos esas mismas combinaciones, siendo $A_1$ cada una de esas intersecciones y $A_2$ nuestro tercer conjunto.

En cada paso tenemos $2^i$ conjuntos disjuntos dos a dos que constituyen una partición de $X$. Cuando lleguemos al final, tendremos la partición que contiene a todos los elementos de ε. 

Es decir, tenemos $N=2^n$ elementos (o menos) disjuntos que generan nuestra $\salgb$. 

En conclusión:
\[\#\algb{M}(ε)=2^N\]
\end{problem}

\begin{problem}[12]
Sellama $\sigma$-anillo de suconjuntos de un conjunto $X$ a toda familia no vacía $\algb{F}$ de subconjuntos de $X$ cerrada para uniones numerables y para las diferencias. Demuestra que todo $\sigma$-anillo es también cerrado para intersecciones numerables. Demuestra que todod $\sigma$-anillo de $X$ es una $\salgb \iff X \in \algb{F}$
\solution
Vamos a probar que un $\sigma$-anillo es cerrado para intersecciones numerables.
Dada una familia infinita numerable $A_i \in \algb{F}$ sabemos que $\bigcup A_i \in \algb{F}$ y lo denotamos por $A$.
Entonces:
\[A \setminus \bigcap A_i = \bigcup(A\setminus A_i) \in \algb{F} \Rightarrow A\setminus (A \setminus \bigcap A_i)=\bigcap A_i \in \algb{F}\]

La condición que le falta a un $\sigma$-anillo para ser una $\salgb$ es contener al total. Por tanto queda clara la última demostración.
\end{problem}

\begin{problem}[13]
Se llama ``clase monótona'' de un conjutno $X$ a toda familia no vacía $\algb{M}$ de subconjuntos de $X$ que sea cerrada para las uniones crecientes y para las intersecciones decrecientes (es decir, si $\forall i=1,2,.. \ C_i \in \algb{M}$ y $C_i \subset C_{i+1}$ o $C_i \supset C_{i+1}$ entonces $\cup_iC_i \in \algb{M}$ o $\cap_iC_i \in \algb{M}$, respectivamente). 

Demuestra que toda $\salgb$ es clase monótona. Da un ejemplo de una clase monótona que no sea $\salgb$
\solution
Antes de nada vamos a definir lo que es una clase monótona.

\begin{defn}[Clase monótona]
$\algb{M} \subset \algb{P}(X)$ es clase monótona si para toda sucesión $\{A_i\}$ de elementos de $\algb{M}$ se cumple que:
\begin{enumerate}
\item \[\forall i A_i \subset A_{i+1} \Rightarrow \bigcup_{n=1}^{\infty}A_i \in \algb{M}\]
\item \[\forall i A_i \supset A_{i+1} \Rightarrow \bigcap_{n=1}^{\infty}A_i \in \algb{M}\]
\end{enumerate}
\end{defn}
Basta con que provemos la primera propiedad ya que la segunda sale por complementación.

No hay nada que demostrar por que esta primera propiedad es parte de la definición de $\salgb$. Es obvio que toda $\salgb$ es clase monótona. La parte interesante del ejercicio es la que sigue.

Vamos a buscar el ejemplo de clase monótona que no sea $\salgb$.
Para el ejemplo basta con encontrar una cadena finita de conjuntos crecientes lo que nos garantiza el cumplimiento de las propiedades de una clase monótona, pero dista mucho de ser una $\salgb$.

Un ejemplo concreto sería:
\[\nat \supset \{2n, \ \forall n \in \nat\}\supset \{4n, \ \forall n \in \nat\}\supset \{8n, \ \forall n \in \nat\}\supset \{16n, \ \forall n \in \nat\}\supset \emptyset\]
La clase monótona sería una subclase de $\algb{P}(\nat)$ formada por todos los conjuntos aquí descritos y no es una $\salgb$, ya que no es cerrada por complementación.
\end{problem}

\begin{problem}[14]
Demuestra que la mínima clase monótona que contiene un álgebra dada $\algb{A}$ es también una $\salgb$

\solution
\textcolor{red}{Atención, ejercicio clave}
\begin{center}
\includegraphics[scale=0.5]{img/clave.jpg}
\end{center}

Vamos a observar dos cosas:
\begin{enumerate}
\item Existe al menos una clase monótona que contiene a un conjunto $\algb{E}$, que es $\algb{P}(X)$
\item La intersección de clases monótonas es clase monótona. Demostración trivial.
\end{enumerate}

Tras estas observaciones podemos hablar de la mínima clase monótona que contiene a $\algb{E}$ como la intersección de todas las clases monótonas que lo contienen y sabemos que la intersección no es vacía por que hay al menos una.

Vamos a tener que hacer dos observaciones más:
\begin{enumerate}
\item Si una clase de conjuntos es álgebra y clase monótona, entonces es $\salgb$.
\item Sea $\algb{A}$ es un álgebra y $\algb{C}$ la mínima clase monótona que contiene a $\algb{A}$. Demostremos que $\algb{C}$ es un álgebra.
\end{enumerate}

%\textcolor{red}{Hecho por mi. No fiarse al 100\%}
% Rehecho en clase y corregidos errores
\begin{proof}
\begin{enumerate}
\item Si nos apoyamos en el ejercicio 6, vemos clara esta demostración. Por ser clase monótona será cerrada para uniones numerables crecientes. Apoyándonos en el ejercicio 6, tenemos que un álgebra cerrada por uniones crecientes numerables es una $\salgb$

\item Vamos a ver que decir que un conjunto $\algb{C}$ es un álgebra es equivalente a decir:
\[\forall E,F \in \algb{C}, E\cap F, \ E\setminus F, \ F\setminus E \in \algb{C}\]

Vamos a definir:
\[\forall E \in \algb{C} \ \algb{C}(E)=\{F \in \algb{C}: \ F \in \algb{C}, E\cap F, \ E\setminus F, \ F\setminus E \in \algb{C}\} \supset \algb{C}\]

Vamos a comprobar que $\algb{C}(E)$ es una clase monótona comprobando las dos propiedades que definen una clase monótona:
\begin{itemize}
\item Dada una sucesión de conjuntos $\{F_i\} \in \algb{C}$ tales que $F_i \subset F_{i+1}$ vemos que:
\ppart \[(\cup F_i)\setminus E = \cup (F_i \setminus E) \in \algb{C}\]
\ppart \[E \setminus (\cup F_i) = \cap(E \setminus F_i) \in \algb{C}\]
\ppart \[E \cap (\cup F_i) = \cup (E \cap F_i) \in \algb{C}\]

De donde podemos deducir que 
\[\cup F_i \in \algb{C} \subset \algb{C}(C)\]

\item Tomemos ahora una sucesión de conjuntos $\{F_i\}\in \algb{C}$ tales que $F_i \supset F_{i+1}$ vemos que:
\ppart \[(\cap F_i)\setminus E = \cap (F_i \setminus E) \in \algb{C}\]
\ppart \[E \setminus (\cap F_i) = \cup(E\setminus F_i) \in \algb{C}\]
\ppart \[E \cap (\cap F_i) = \cap (E \cap F_i)\in \algb{C}\]

De donde podemos deducir que, en este caso:
\[\cap F_i \in \algb{C} \subset \algb{C}(C)\]
%TODO el final de cada parte no me convence
\end{itemize} 

Además esta claro que:
\[\forall E,F \in \algb{C} \implies E\in \algb{C}(F) \iff F \algb{C}(E)\]

Vamos a ver ahora que $\algb{C} \subset \algb{C}(E) \forall E \in \algb{A}$

Para probarlo, basta con ver que $\algb{A} \in \algb{C}(E)$ ya que $\algb{C}$ es la mínima clase monótona que contiene a $\algb{A}$.

Pero por lo visto anteriormente, tenemos que:
\[\forall E \in \algb{A}, \ \forall F \in \algb{C} \Rightarrow E \in \algb{C}(F) \Rightarrow \algb{A} \subset \algb{C}(F) \Rightarrow \algb{C} \subset \algb{C}(F)\]

\end{enumerate}
\end{proof}
\end{problem}

\subsection{Hoja 3}
A lo largo de esta hoja se piden realizar varias demostraciones. La terminología que ha establecido Patricio en estos ejercicios consiste en llamar comprobaciones a los ejercicios mas triviales, pruebas a los del un nivel intermedio y demostraciones a los que realmente requieren esfuerzo.

\begin{problem}
Sea $(X, \algbM, µ)$ un espacio de medida. Si $E,F \in \algbM$ comprueba que 
\[µ(E)+µ(F)=µ(E\cup F) - µ (E\cap F)\]

\solution
En el enunciado pide comprobar por lo que, según Patricio, debe ser bastante trivial.

En este caso tiene razón, ya que lo ocurre es conceptualmente muy sencillo. Al sumar $µ(E)+µ(F)$ la intersección la estoy teniendo en cuenta dos veces (una al medir $E$ y otra al medir $F$).

Para escribir una demostración formal observemos que:
\[E \cup F = (E \setminus F) \cup (F\setminus E) \cup (E \cap F)\]
Puesto que se trata de una unión disjunta, y µ es una medida, si tomamos la medida de todo ello obtenemos:
\[µ(E \cup F) = µ(E \setminus F) + µ(F\setminus E) + µ(E \cap F) = µ(E)-µ(F\cap E)+µ(F)-µ(E \cap F)+µ(E \cap F) =\]
\[= µ(E) + µ (F) +µ (E \cap F)\]
\end{problem}

\begin{problem}
Sea $(X, \algbM, µ)$ un espacio de medida y sea $E \in \algbM$. Para cada $A \in \algbM$ sea $µ_E(A)=µ(E\cap A)$. 

Comprueba que $µ_E$ es una medida en $\algbM$
\solution
Si recordamos la definición de medida en un espacio medible, vemos que se trata de una aplicación de los subconjuntos del espacio a los reales que cumple dos propiedades básicas:
\begin{enumerate}
\item La medida del vacío es 0
\item La medida de una unión de conjuntos disjuntos es la suma de las medidas de cada uno de los conjuntos
\end{enumerate}

Si queremos probar que $µ_E$ es una medida deberemos comprobar que se satisfacen esas dos conciciones, es decir:
\[µ_E(\emptyset) = 0\]
\[µ_E(\bigcup A_n) = \sum µ_E(A_n)\]

La primera propiedad es trivial, ya que, por ser µ una medida tenemos:
\[µ_E(\emptyset)=µ(E \cap \emptyset)=µ(\emptyset)=0\]

Para la segunda propiedad vemos que:
\[µ_E(\bigcup A_n)=µ_E(E \cap (\bigcup A_n)) = µ_E(\bigcup(A_n \cap E)) \text{ con } A_i\cap A_j = \emptyset \ \forall i\neq j\]
Puesto que los $A_n \cap E$ son disjuntos y $µ$ es una medida sobre $\algbM$ tenemos que:
\[µ(\bigcup(A_n \cap E))=\sum µ(E \cap A_n) = \sum µ_E(A_n)\]

\end{problem}
\newpage
\begin{problem}
\begin{enumerate}
\item Comprueba que una medida $\sfin$ es semifinita
\item Sea $X$ un conjunto no numerable y sea µ la medida discreta en $(X, \algbP (X))$, comprueba que µ es  semifinita pero no $\sfin$
\end{enumerate}

\solution
\begin{enumerate}
\item
Sabemos que, por definición de $\sfin$, el conjunto $X$ (sobre el que está definida la medida) puede expresarse como una unión de elementos de la $\salgb$ de medida finita. Es decir:
\[\exists \{E_n\}_{n \in \nat} \subset \algbM \tq X = \bigcup E_n \text{ con } µ(E_n)< \infty \forall n\]

Y para ver que µ es semifinita, por definición, debemos ver que todo subconjunto $E$ de la $\salgb$ tiene un subconjunto contenido también en la $\salgb$ de medida finita. Es decir:
\[\forall E \subset \algbM \text{ con } µ(E)>0, \ \exists A \in \algbM \text{ con } A\neq \emptyset \ A \subset E \text{ y } 0<µ(A) < \infty\]

Puesto que $X$ puede expresarse como unión de los $E_n$ está claro que:
\[E= \bigcup (E \cap E_n)\]
y puesto que no podemos garantizar que sean disjuntos los elementos de la unión tenemos:
\[µ(E)=µ(\bigcup (E \cap E_n))\leq \sum(µ(E \cap E_n))\]

Entonces 
%TODO WTF? 
\[\exists n \tq 0<µ(E \cap E_n) < \infty\]
Y basta con tomar $A=E \cap E_n \subset E$

\item 
Recordemos que la medida discreta es la que, aplicada sobre un conjunto, nos devuelve el cardinal del conjunto si este es finito o inftinito en caso contrario.

Para ver que es semifinita nos fijamos en que
\[µ(E) > 0 \implies E \neq \emptyset \implies \exists a \in E\]
En ese caso 
\[\{a\}\subset E \text{ con } µ(\{a\})=1\]
Por lo que es semifinita 
%TODO WTF?

Ahora vamos a comprobar que no es $\sfin$ por reducción al absurdo.

Si fuese $\sfin$,
\[\exists \{E_n\}_{n \in \nat} \subset \algbM \tq X = \bigcup E_n \text{ con } µ(E_n)< \infty \forall n\]
es decir, el cardinal de $E_n$ sería finito  y $X= \bigcup_{n \in \nat }E_n$ lo que implicaría que $X$ es numerable, en contradición la hipótesis inicial.

\end{enumerate}
\end{problem}

\begin{problem}
Sea $X$ un conjunto finito y numerable, para $A \subset X$ se define
$µ(A)=0$ si A es finito y $µ(A)=\infty$ si A infinito.

\begin{enumerate}
\item Comprueba que µ es finitamente aditiva (en $(X, \algbP (X))$) pero no numerable aditiva
\item Comprueba que existe una sucesión creciente de conjuntos $\{A_n\}$ tal que:
\[\forall n \in \nat, \ µ(A_n)=0 \ y \ \lim_{n} A_n =X\]
\end{enumerate}
\solution
\begin{enumerate}
\item Si tomamos una sucesión finita y disjunta de $A_n$, si todos son finitos las medidas serán 0, la unión será finita y también tendrá medida 0.

Si al menos uno es infinito, entonces la medida sería infinita y la suma de medidas será infinito también.

Sin embargo, si tomo una sucesión numerable de conjuntos finitos, la medida de cada uno de ellos sería 0 pero la unión sería infinita por lo que tendría medida infinita.

Para obtener elegamentemente estos conjuntos podemos tomar una numeración de los elementos de X (ya que es numerable). Es decir, consideramos:
\[X= \bigcup_{n=0}^{\infty}A_n \text{ donde } A_n=\{a_n\}\subset X\]
En este caso µ(X) es finita pero a la izquierda tendríamos una suma infinita de valores positivos ($µ(A_n)$) que nos daría infinito

\end{enumerate}

\end{problem}

\begin{problem}
Sea ($X, \algbM, µ$) un espacio de medida y $\{E_n\} $ una sucesión de conjuntos en $\algbM$. Demuestra que si existe un k tal que $µ(\bigcup_{i=k}^{\infty}E_i) < \infty$ entonces:
\[µ(\liminf (E_j))<\liminf (µ(E_j)))\]
\[µ(\limsup (E_j))>\limsup (µ(E_j)))\]

En particular si µ(X) < $\infty$ entonces:
\begin{enumerate}
\item 
\[µ(\liminf E_j) \leq \liminf µ(E_j) \leq \limsup µ(E_j) \leq µ(\limsup(E_j))\]
\item Si existe $\lim E_j$ entonces $µ(\lim E_j)=\lim µ(E_j)$.
\end{enumerate}
Indica en qué punto es necesaria la condición $µ(\bigcup_{j=k}^{\infty}E_j)< \infty$ para al menos un k.
\solution
Recordemos algunas 'definiciones' necesarias para resolver el ejercicio:
\[\liminf E_k)=\bigcup_{j=0}^{\infty}\bigcap_{k=j}^{\infty}E_k\]
\[\limsup E_k)=\bigcap_{j=0}^{\infty}\bigcup_{k=j}^{\infty}E_k\]
Basándonos en la Wikipedia, daremos una definición más intuitiva de estos conceptos:

\begin{defn}[Límite superior]
Es el conjunto formado por los elementos que pertenecen a todos los conjuntos de la sucesión salvo, quizás, a un número finito de ellos
\end{defn}

\begin{defn}[Límite superior]
Es el conjunto formado por todos los elementos que pertenecen a infinitos conjuntos en la sucesión
\end{defn}

Además, como ya vimos, si una sucesión es creciente el límite coincide con el límite de la unión. Igual ocurre con una sucesión decreciente y la intersección. Es decir:
\[E_i \subset E_{i+1} \implies \lim E_i = \bigcup E_i\]
\[E_i \supset E_{i+1} \implies \lim E_i = \bigcap E_i\]


Vamos ya a por el ejercicio.

Si tomamos un elemento $E_k$ y lo intersecamos con uno o varios elementos de la sucesión obtendremos un subconjunto de $E_k$, es decir:
\[E_k \supset \bigcap_{j=k}^{\infty}E_j\]
Por tanto, si aplicamos medidas a ambos lados tenemos:
\[µ(E_k) < µ(\bigcap_{j=k}^{\infty}E_j) \rightarrow µ(\bigcup_{k=0}^{\infty}\bigcap_{j=k}^{\infty}E_j)\]

La convergencia se debe a que la intersección será más grande (estoy tomando menos conjuntos) por lo que al final me quedarán aquellos elementos que estén en todos los conjuntos salvo en un número finito de ellos. (Está en todos los conjuntos salvo, quizás, en los que he ido quitando)

De esto, y basándonos en las definiciones iniciales, podemos concluir:
\[\lim \inf µ(E_k) \geq \lim \inf µ(\bigcap_{j=k}^{\infty}E_j) = \lim (µ(\bigcap_{j=k}^{\infty}E_k))=µ(\bigcup_{k=0}^{\infty}\bigcap_{j=k}^{\infty}E_j)\]
%TODO No veo claro el por qué

Ahora bien, tenemos una sucesión de conjuntos de la que caculamos el límite superior y el límite inferior mediante las dos definiciones proporcionadas al inicio del ejercicio.

Si tiramos los n primeros conjuntos y nos quedamos con los demás, el límite superior y el límite inferior siguen siendo los mismos.

Analicemos un poco más en detalle esta afirmación:
\[x \in \lim \sup (E_i)_{i=0}^{\infty} \text{ si } \forall k \exists k' \geq k \tq x \in E_{k'}\]
%TODO Ha copiado otro par de definiciones que yo no tengo. apañarlo.

Una vez queda claro que el límite superior y el inferior no dependen de los n primeros elementos de la familia, está claro que, sin pérdida de generalidad, podemos demotrar lo que pide el enunciado suponiendo que k=0.

Tenemos pues:
\[µ(\bigcup_{i=0}^{\infty} E_i)< \infty\]
Basándonos en las definiciones iniciales, vemos que lo que nos pide demostrar el enunciado es equivalente a:
\[µ(\bigcap_{k=0}^{\infty}\bigcup_{j=k}^{\infty}E_j)= µ(E)-µ((\bigcap_{k=0}^{\infty}\bigcup_{j=k}^{\infty}E_j)^c)\]
siendo $E=\bigcup_{j=0}^{\infty}E_j$ y hablando del complementario como el complementario dentro de $E$

Pero si lo analizamos vemos que lo que tenemos es:
\[\limsup µ(\bigcup_{k=0}^{\infty}\bigcap_{j=k}^{\infty}E_j)=\]
\[=µ(E)-\liminf µ(\bigcup_{k=0}^{\infty}\bigcap_{j=k}^{\infty}E^c_j) \geq \]
\[\geq µ(E)-\liminf(µ(E)-µ(E_k))=µ(E)-µ(E)+\limsup µ(E_k)\]

Tomando el incio y el final de esta cadena de desigualdades llegamos a lo que queríamos demostrar:
\[\limsup µ(\bigcup_{k=0}^{\infty}\bigcap_{j=k}^{\infty}E_j)=\limsup µ(E_k)\]
\end{problem}

\begin{problem}
Sea $X=\{a_1, a_2, a_3\}$ y µ una medida definida en $\algbP (X)$ tal que $µ(A_i)=\frac{1}{3} \ \forall i$.

Se define una sucesión $A_n$ tal que $A_{2k}=\{a_1, a_2\}$ y $A_{2k+1}=\{a_3\}$ $\forall k \in \nat$.

Prueba que:
%TODO cadena de desigualdades del ejercicio anterior
\solution
En este caso está bastante claro que:
\[\liminf A_n = \emptyset\]
puesto que este límite consiste en tomar la unión de todos los elementos que están presentes en todos los elementos de la sucesión. No hay ningún elemento que esté en todos los elementos y por ello obtenemos el vacío.
Así mismo
\[\limsup A_n = X\]
ya que este límite implica tomar la intersección de conjuntos formados por los elementos contenidos en algún elemento de la sucesión. Todos los elementos de $X$ están contenidos en algún elemento de la sucesión y la intersección de los $X$ con ellos mismo es $X$.

Por otro lado:
\[\limsup µ(A_n) = \frac{2}{3}\]
\[\liminf µ(A_n) = \frac{1}{3}\]

Puesto que $µ(X) =1$ y $µ(\emptyset)=0$ obtenemos que es claramente correcta la cadena de desigualdades.
\end{problem}


\begin{problem}
Sea µ una medida semifinta y sea $E$ tal que $µ(E) < \infty$. Prueba que si c es un número real mayor que 0, existe un conjunto $F \subset E$ tal que c < $µ(F)$ < $\infty$

\solution 
Basándonos en la sugerencia del enunciado, tomemos:
\[k=\sup \{µ(F): \ F \subset E, \ µ(F)< \infty \}\]

Por ser k un supremo,
\[\forall n \exists F_n \subset E \tq k-\frac{1}{n} \leq µ(F_n) \leq k\]
pero
\[k - \frac{1}{n} \geq µ(\bigcup_{n-1}^{N}) \leq k\]
por lo que si tomamos un n suficientemente grande llegaremos a la igualdad:
\[k \leq µ(\bigcup_{n=1}^{\infty}F_n) \leq k\]
Ahora podemos construir una sucesión creciente de $F_n$ tales que la unión de todos ellos es F. Es decir:
\[F= \bigcup F_n\]
y como la medida de $F$ es finita sabemos que:
\[µ(E \setminus F)= \infty\]
y por tanto
\[\exists F'\subset (E \setminus F) \tq µ(F')> 0\]

Entonces
\[F \cup F' \subset E \Rightarrow µ(F \cup F')=µ(F)+µ(F') > k\]
Y llegamos a una contradicción porque k era el supremo.


\end{problem}

\begin{problem}
En un espacio de medida $(X, \algbM, µ)$ sea $\{A_i\}$ una familia de conjuntos medibles tales que $\sum µ(A_i) < \infty$.

Prueba que casi todo elemento $x \in X$ pertenece sólo a un número finito de $A_i$

\solution

Recordemos que el conjunto de los $x \in X$ que pertenecían a un número infinito de $A_i \subset X$ lo denominamos límite superior.

\[\limsup = \bigcap_{n=1}^{\infty}\bigcup_{k=n}^{\infty} A_k\] 

Vemos ahora que pasa con su medida:
\[µ(\limsup) = µ(\bigcap_{n=1}^{\infty}\bigcup_{k=n}^{\infty} A_k) \leq \sum_{k=n}^{\infty}µ(A_n) \rightarrow 0\]

La última convergencia aquí indicada se deduce de la finitud de la serie. Si una serie es finita, su cola converge a 0.
\end{problem}

\begin{problem}
Sea $X$ un conjunto no numerable. Sea $\algbM$ la $\salgb$ formada por los conjuntos finitos o numerables y los conjuntos con complementario finito o numerable. Sea µ la medida que me da el cardinal del conjunto si éste es finito e infinito si el conjunto es infinito:
\begin{enumerate}
\item Demuestra que µ es una medida completa en $\algbM$
\item Estudia la medida $µ^*$ construida a partir de µ y de $\algbM$
\end{enumerate} 

\solution
Aunque el enunciado no dice mucho al respecto, Patricio insiste en dejar claro que  $\algbM$ es realmente una $\salgb$ lo cual se ve de manera sencilla.

\begin{enumerate}
\item Para ver que se una medida completa debemos ver que la medida del vacío es 0, cosa que es obvia pues el cardinal del vacío es 0.

Además debemos comprobar que todo subconjunto de un conjunto de medida 0 se contiene en $\algbM$.

Puesto que en $\algbM$ el único conjunto de medida 0 es $\emptyset$, no hay nada que demostrar.

\item La medida exterior a partir de µ y de $\algbM$.

En este caso, $µ^*(E)=card(E)$ si $E$ es finito y $µ^*(E)=\infty$ en caso contrario.

Habría que ver ahora qué ocurre con los conjuntos $\notin \algbM$. Para ellos todo recubrimiento contiene algún conjunto ... bla bla bla
%TODO...
\end{enumerate}
\end{problem}

\begin{problem}
Se define $µ^*$ sobre $\algbP (\nat)$ como:
\[µ^*(E) = \frac{n}{n+1} \text{ si } E \text{ es finito }\]
\[µ^*(E) = \infty \text{ si } E \text{ es infinito}\]
Demuestra que $µ^*$ es una medida exterior y halla la $\salgb$ de los conjuntos medibles
\solution

Para ver que $µ^*$ es una medida exterior debemos que:
\begin{itemize}
\item $µ^*(\emptyset)= 0$

\item $E \subset E' \implies µ^*(E)<µ^*(E')$

Vemos que en todos los casos posibles para estos $E$ se verifica la desigualdad: los dos finitos, infinitos, uno y uno
%TODO completar

\item $µ^*(\bigcup A_n) \leq \sum µ^*(A_n)$
Como en el apartado anterior, podemos verlo por casos.

%TODO por casos otra vez, exactamente como lo hice en casa
\end{itemize}

Vamos a ver ahora cuál es la $\salgb$ de los conjuntos medibles.
\[\algbM^*=\{E \subset \nat \ \forall A \subset \nat \ µ^*(A)\geq µ^*(A \cap E) + µ^*(A \cap E^c)\}\]
Tomamos la desigualdad porque la desigualdad contraria está garantizada siempre, de modo que forzar que se cumpla esta desigualdad implica el cumplimiento de la igualdad.

Para que un conjunto $E$ pertenece a esta $\salgb$ necesitamos ver que la parte de la derecha de la desigualdad es menor que 1.

Sin embargo, tanto tomando $E$ finito como infinito nos tomamos siempre con contradicción. Al final este algebra solo tiene el vacío y el total.
\end{problem}

\begin{problem}
Dada una medida exterior $µ^*$ en $X$ y $\{A_j\}_{j \in \nat}$ una familia disjunta de conjuntos $µ^*-$medibles. 

Demuestra que para cualquier $E \subset X$, se cumple que:
\[µ^*(E \bigcap (\bigcup_{j=0}^{\infty} A_i))= \sum_{j=0}^{\infty}µ^*(E \bigcap A_j)\]
\solution
caca
\end{problem}